\section{Character Sums via the Quadratic Character}
For certain special characters, the associated Gaussian sums can be evaluated explicitly. A very special multiplicative character to study is the quadratic character, $\eta$. We get a very celebrated formula for the Gaussian sum of $\eta$ and the canonical additive character. 
\begin{Theorem}{}{}
  Let $\bbF_q$ be a finite field with $q=p^r$ where $p$ is an odd prime and $r\in\bbN$. Let $\eta$ be the quadratic character of $\bbF_q$ and let $\chi_1$ be the canonical additive character of $\bbF_q$. Then $$G(\eta,\chi_1)=\begin{cases}
    (-1)^{r-1}\cdot q^{\nicefrac12} & \text{if }p\equiv 1\bmod 4\\ 
    (-1)^{r-1}\cdot i^{r}\cdot q^{\nicefrac12} & \text{if }p\equiv 3\bmod 4    
  \end{cases}$$
  In other words, $$G(\eta,\chi_1)=(-1)^{r-1}\cdot i^{\frac{(p-1)^2}4 r} \cdot q^{\nicefrac12}$$ 
\end{Theorem}
\begin{proof}
  Since $\eta$ always takes values in $\pm 1$ we have $\eta=\ov{\eta}$. Hence by using \lemref[(iv)]{lem:gaussian-sum-identity} we have $G(\eta,\chi_1)^2=\eta(-1)q$. Now if $q\equiv 1\bmod 4$ then $\eta(-1)=1$ and if $q\equiv 3\bmod 4$ then $\eta(-1)=-1$. Hence \begin{equation}
    G(\eta,\chi_1)=\begin{cases}
      \pm\ q^{\nicefrac12}& \text{if }q\equiv 1\bmod 4\\ 
      \pm\ i\cdot q^{\nicefrac12}& \text{if }q\equiv 3\bmod 4
    \end{cases}
  \end{equation}
  Since $i^{\nicefrac{(p-1)^2}{4}}=\begin{cases}
      1 & \text{if }q\equiv 1\bmod 4\\ 0 & \text{if }q\equiv 3\bmod 4
    \end{cases}$

  \begin{observation}\label{obs:gauss-quad-char}
    For any finite field $\bbF_q$,  $G(\eta,\chi_1)=\pm\ i^{\nicefrac{(q-1)^2}{4}}\cdot q^{\nicefrac12}$ 
  \end{observation}

  Now comes the hardest part of the proof that is getting the signs correctly. We will first prove this for $r=1$ case then we will go for $r>1$. Now we have the observation
  
  
  Let $W$ be the vector space of all complex valued functions on $\bbF_p^*$. It is a $(p-1)$-dimensional vector space over $\bbC$. So $\mcB\coloneqq \{f_1,\dots, f_{p-1}\}$ where $f_j(k)=1$ if $j=k$ and other wise $0$ forms a basis of $W$ for all $j\in[p-1]$. Now by using the orthogonal properties of multiplicative characters in \thmref{thm:mult-char-props} the set $\mcB_{\Mq}\coloneqq \{\psi_0,\dots, \psi_{p-2}\}$ as defined in \lemref{lem:mult-char-define}, also forms a basis of $W$.

  Let $\zeta_p$ is the $p^{th}$ root of unity i.e. $\zeta_p=e^{\nicefrac{2\pi i}{p}}$. Hence if $\chi_k\in\Xp$ be an additive character of $\bbF_p$ as defined in \thmref{thm:add-char-define}, then for all $j\in\bbF_p$, $\chi_k(j)=\zeta_p^{jk}$. Then define the linear operator $T:W\to W$ in the following way:
  \begin{equation}\label{eq:quad-linear-operator}
    \forall\ h\in W,\ \forall\ k\in[p-1]\qquad (Th)(k)=\sum_{j=1}^{p-1}\zeta_p^{jk}h(j)
  \end{equation}
  Therefore for all basis elements $\psi_i$ $$(T\psi_i)(k)=\sum_{j=1}^{p-1}\zeta_p^{jk}\cdot \psi_i(k)=\sum_{j=1}^{p-1}\chi_k(j)\cdot \psi_i(k)=G(\psi,\chi_k)=\ov{\psi_i}(k)\cdot G(\psi_i,\chi_1)\quad \forall\ k\in\bbF_p^*$$where the last equality comes from \lemref[(i)]{lem:gaussian-sum-identity}. Hence $T\psi_i=G(\psi_i,\chi_1)\cdot \ov{\psi_i}$ for every multiplicative characters of $\bbF_p$. 

  Since $\ov{\psi}=\psi$ only for the trivial character and quadratic character, the matrix for $T$ in the basis $\mcB_{\Mq}=\{\psi_0,\dots, \psi_{p-2}\}$ contains two diagonal entries, $G(\psi_0,\chi_1)=-1$ and $G(\eta,\chi_1)$ for $\psi_0$ and $\eta=\psi_{\nicefrac{p-1}2}$. For all other characters we can divide them into groups of pairs where the characters in each pair are conjugates of each other. Then for each pair $(\psi,\ov{\psi})$ we have the following block matrix in the matrix of $T$. $$\mat{0 & G(\ov{\psi},\chi_1)\\ G(\psi,\chi_1)& 0}$$Therefore in the determinant of the matrix $T$ each block contributes $-G(\psi,\chi_1)G(\ov{\psi},\chi_1)=-\psi(-1)p$ by \lemref[(iv)]{lem:gaussian-sum-identity}. So we get $$\det(T)=(-1)\cdot G(\eta,\chi_1) \cdot (-p)^{\nicefrac{(p-3)}2}\prod_{j=1}^{\frac{p-3}2}\psi_j(-1)$$Let $\psi\in\Mp$ is the principle multiplicative character of $\bbF_p$. Then $\psi_j=\psi^j$. Therefore $\psi_j(-1)=\psi^j(-1)=(-1)^j$ as $\psi(-1)=-1$. Hence we have \begin{align*}
    \det(T) & = - G(\eta,\chi_1) \cdot (-p)^{\nicefrac{(p-3)}2}\prod_{j=1}^{\frac{p-3}2}(-1)^j\\ 
    & = - G(\eta,\chi_1) \cdot (-p)^{\nicefrac{(p-3)}2}\cdot (-1)^{\sum\limits_{j=1}^{\nicefrac{(p-3)}{2}}j}\\ 
    & =  (-1)^{\nicefrac{(p-1)}{2}}\cdot G(\eta,\chi_1) \cdot p^{\nicefrac{(p-3)}2}\cdot (-1)^{\nicefrac{(p-1)(p-3)}8}\\ 
    & = (-1)^{\nicefrac{(p-1)}{2}}\cdot \lt(\pm\ i^{\nicefrac{(p-1)^2)}{4}}\cdot p^{\nicefrac12}\rt) \cdot p^{\nicefrac{(p-3)}2}\cdot (-1)^{\nicefrac{(p-1)(p-3)}8}\by{\obsref{obs:gauss-quad-char}}\\ 
    & = \pm (-1)^{\nicefrac{(p-1)}{2}}\cdot i^{\frac{(p-1)^2)}{4}+\frac{(p-1)(p-3)}4}\cdot p^{\nicefrac{(p-2)}{2}}\\ 
    & = \pm (-1)^{\nicefrac{(p-1)}{2}}\cdot i^{\nicefrac{(p-1)(p-2)}2}\cdot p^{\nicefrac{(p-2)}{2}}
  \end{align*}
  So we get \begin{equation}\label{eq:fourier-basis-det}
    \det(T)=\pm (-1)^{\nicefrac{(p-1)}{2}}\cdot i^{\nicefrac{(p-1)(p-2)}2}\cdot p^{\nicefrac{(p-2)}{2}}
  \end{equation}

  Now we will compute $\det(T)$ utilizing the basis $\mcB=\{f_1,\dots, f_{p-1}\}$. From the construction of $T$ in \eqref{eq:quad-linear-operator} we get the $(j,k)^{th}$ element of the matrix of $T$ in the $\{f_1,\dots, f_{p-1}\}$ basis to be $\zeta_p^{jk}$ for any $0\leq j,k\leq p-1$. So we have \begin{align*}
    \det(T)& =\det\lt( \lt\{\zeta_p^{jk}\rt\}_{1\leq j,k\leq p-1}\rt)\\ 
    & = \det\lt( \lt\{\zeta_p^j\zeta_p^{j(k-1)}\rt\}_{1\leq j,k\leq p-1}\rt)\\ 
    & = \zeta_p^{1+2+\cdots+(p-1)}\det\lt( \lt\{\zeta_p^{j(k-1)}\rt\}_{1\leq j,k\leq p-1}\rt)\\ 
    & = \det\lt( \lt\{\zeta_p^{j(k-1)}\rt\}_{1\leq j,k\leq p-1}\rt) & \lt[\text{As }p\mid \frac{p(p-1)}2=1+2+\cdots+(p-1)\rt]
  \end{align*}
  Now the matrix $\lt\{\zeta_p^{j(k-1)}\rt\}_{1\leq j,k\leq p-1}$ is the \emph{Vandermonde} matrix $V(\zeta_p,\zeta_p^2,\dots, \zeta_p^{p-1})$. Therefore $$\det(T)=\prod_{1\leq j<k\leq p-1}\lt(\zeta_p^k-\zeta_p^j\rt)$$We will show that 
  \begin{equation}\label{eq:quadratic-gaussian-det}
    (-1)^{\nicefrac{(p-1)}{2}}\cdot i^{\nicefrac{(p-1)(p-2)}2}\cdot p^{\nicefrac{(p-2)}{2}}=\prod_{1\leq j<k\leq p-1}\lt(\zeta_p^k-\zeta_p^j\rt)
  \end{equation}
  So let $\zeta_{2p}=e^{\nicefrac{\pi i}{p}}$. Hence $\zeta_{2p}^2=\zeta_p$. By the properties of complex numbers we know for all $t\in\bbZ$, $\zeta_{2p}^t=\cos\lt(\frac{t\pi}{p}\rt)+i\sin\lt(\frac{t\pi}p\rt)$. Hence $\zeta_{2p}^t-\zeta_{2p}^{-t}=2i\sin\lt(\frac{t\pi}p\rt)$. So we get \begin{align*}
    \det(T) & = \prod_{1\leq j<k\leq p-1}\lt(\zeta_p^k-\zeta_p^j\rt)\\
    & = \prod_{1\leq j<k\leq p-1}\lt(\zeta_{2p}^{2k}-\zeta_{2p}^{2j}\rt)\\
    & = \prod_{1\leq j<k\leq p-1}\zeta_{2p}^{j+k}\lt(\zeta_{2p}^{k-j}-\zeta_{2p}^{-(k-j)}\rt)\\ 
    & =  \lt[\prod_{1\leq j<k\leq p-1}\zeta_{2p}^{j+k} \rt]\;\lt[\prod_{1\leq j<k\leq p-1} \lt(2i\sin\frac{\pi(k-j)}{p}\rt) \rt]\\ 
    & = \zeta_{2p}^{\sum_{1\leq j<k\leq p-1}(j+k)} \cdot i\st^{\binom{p-1}{2}}\prod_{1\leq j<k\leq p-1} \lt(2\sin\frac{\pi(k-j)}{p}\rt) \\ 
    & = \zeta_{2p}^{\sum_{1\leq j<k\leq p-1}(j+k)} \cdot i{\nicefrac{(p-1)(p-2)}  {2}}\prod_{1\leq j<k\leq p-1} \lt(2\sin\frac{\pi(k-j)}{p}\rt) 
  \end{align*}
  Now we calculate the power of $\zeta_{2p}$ in the first term. $$\sum_{1\leq j<k\leq p-1}(j+k)=\sum_{k=2}^{p-1}\sum_{j=1}^{k-1}(j+k)=\sum_{k=2}^{p-1}\lt(\frac{k(k-1)}2+k(k-1)\rt)=\frac32\sum_{k=2}^{p-1}(k^2+k)=\frac{p(p-1)(p-2)}2$$Hence $$\zeta_{2p}^{\sum_{1\leq j<k\leq p-1}(j+k)} = \lt(\zeta_{2p}^p\rt)^{\nicefrac{(p-1)(p-2)}{2}}=(-1)^{\nicefrac{(p-1)(p-2)}{2}}=\lt((-1)^{p-2}\rt)^{\nicefrac{(p-1)}{2}}=(-1)^{\nicefrac{(p-1)}{2}}$$So we have \begin{equation}\label{eq:standard-basis-det}
    \det(T)=(-1)^{\nicefrac{(p-1)}{2}}\cdot i^{\nicefrac{(p-1)(p-2)}{2}}\cdot \prod_{1\leq j<k\leq p-1} \lt(2\sin\frac{\pi(k-j)}{p}\rt)
  \end{equation}Notice if we compare this with \eqref{eq:fourier-basis-det}, the first two terms are already same. If we can show that $\prod_{1\leq j<k\leq p-1} 2\sin\frac{\pi(k-j)}{p}>0$ then we have \eqref{eq:quadratic-gaussian-det}. 
  
  Now since $j,k\in[p-1]$ and $k>j$ we have $k-j\in[p-1]$. Hence $\sin\frac{\pi(k-j)}{p}>0$ for all $1\leq j<k\leq p-1$. Hence $$\prod_{1\leq j<k\leq p-1} \lt(2\sin\frac{\pi(k-j)}{p}\rt) >0$$ So if we compare \eqref{eq:standard-basis-det} with \eqref{eq:fourier-basis-det} and match the signs we get $$\det(T)=(-1)^{\nicefrac{(p-1)}{2}}\cdot i^{\nicefrac{(p-1)(p-2)}2}\cdot p^{\nicefrac{(p-2)}{2}}$$So from this we get $G(\eta,\chi_1)=i^{\nicefrac{(p-1)^2}{4}}\cdot p^{\nicefrac12}$ since at no point of our calculation the sign was changed. 

  So now we have established the theorem for $r=1$. Now for $r>1$ following \autoref{sec:lifting-chars} we can lift the quadratic character of $\bbF_p$, $\eta$ to a multiplicative character $\eta^{(r)}$ of $\bbF_q$ where $q=p^r$ and $\eta^{(r)}=\eta\circ N$ where $N$ is the absolute norm function from $\bbF_q$ to $\bbF_p$. Since $\eta$ is a quadratic character $\eta^{(r)}$ will be the quadratic character of $\bbF_q$. Similarly we can lift the canonical additive character $\chi_1$ of $\bbF_p$ to the canonical additive character $\chi_q^{(r)}$ of $\bbF_q$ by $\chi_1^{(r)}=\chi_1\circ \tr$ where $\tr$ is the absolute trace function from $\bbF_q$ to $\bbF_p$. Then by using \DHthm we get $$    G(\eta^{(r)},\chi_1^{(r)}) = (-1)^{r-1}G(\eta,\chi_1)^r = (-1)^{r-1}\lt[i^{\nicefrac{(p-1)^2}{4}}\cdot p^{\nicefrac12}\rt]=(-1)^r \cdot i^{\frac{(p-1)^2}{4}r}\cdot p^{\nicefrac{r}2}=(-1)^r \cdot i^{\frac{(p-1)^2}{4}r}\cdot q^{\nicefrac12}$$So we have the theorem.
\end{proof}

The proof of the above theorem gives us another important result. 
\begin{lemma}{}{}
  If $p$ is a odd prime and $\zeta_p$ is the $p^{th}$ root of unity. Let $V(\zeta_p,\zeta_p^2,\dots,\zeta_p^{p-1})$ is the $p\times p$ Vandermonde Matrix. Then $$\det\lt(V(\zeta_p,\zeta_p^2,\dots,\zeta_p^{p-1})\rt)=(-1)^{\nicefrac{(p-1)}{2}}\cdot i^{\nicefrac{(p-1)(p-2)}2}\cdot p^{\nicefrac{(p-2)}{2}}$$
\end{lemma}